% ================================================================
% ==== FILE: content/m_spaces/m6.tex
% ================================================================

% ==============================
%  Ontology of Continua — Core
%  Meta-Space: M6
%  FULL MODULE — FINAL VERSION
% ==============================

\subsubsection{\texorpdfstring{$M_6$}{M_6} Übersicht}
\label{sec:m6-overview}

$M_6$ ist der Metaraum, der \emph{kognitive Organisation} zulässig macht.
Während $M_5$ elektrische Erregbarkeit und Spitzenzyklen unterstützt, ermöglicht $M_6$:

\begin{itemize}
    \item stabile Attraktoren über Spike-Zügen,
    \item semantische Achsen und Darstellungsgeometrie,
    \item protologische Operationen, die auf neuronale Zustände wirken,
    \item-Klassifizierung und Vorhersagedynamik,
    \item Entstehung symbolischer und subsymbolischer Konzepte,
    \item Kohärenzbeschränkungen, die für kognitive Kontinua erforderlich sind ($K_6$).
\end{itemize}

Somit ist $M_6$ die minimale Umgebung für \emph{bedeutungstragende neuronale Dynamik}.

% ---------------------------------------------------------------
\subsubsection*{1. Zulässiger Konfigurationsraum \texorpdfstring{$\Omega(M_6)$}{\Omega(M_6)}}

\[
\Omega(M_6) =
\left\{
(S(t),\ A_{\mathrm{sem}},\
\mathcal{W},\
\mathcal{A},\
C_{\mathrm{att}},\
\Pi_6)
\ \big|\ \text{Zulässigkeitsbedingungen gelten}
\right\}.
\]

Wo:

\begin{itemize}
    \item $S(t)$ – neuronale Zustandsbahnen (Spike Trains, Feuerraten, Phasencodes),
    \item $\mathcal{W}$ – synaptischer Gewichtstensor mit Plastizitätsbeschränkungen,
    \item $\mathcal{A}$ — zulässiger Attraktorsatz,
    \item $A_{\mathrm{sem}}$ – semantische Achsen eingebettet in die Attraktorgeometrie,
    \item $C_{\mathrm{att}}$ — Zyklusfamilie von Aufmerksamkeit / Arbeitsgedächtnis,
    \item $\Pi_6$ – Einschränkungen, die Stabilität und Trennbarkeit von Darstellungen gewährleisten.
\end{itemize}

Eine Konfiguration ist zulässig, wenn:

\begin{enumerate}
    \item neuronale Dynamik kann zu mindestens einem stabilen Attraktor konvergieren;
    \item Attraktorbecken sind durch Hyperebenen oder Mannigfaltigkeiten trennbar;
    Die \item-Zuordnung zwischen Eingabemustern und Attraktoren ist konsistent.
    \item-Plastizitätsregeln bewahren die Kohärenz, anstatt Strukturen zu zerstören.
\end{enumerate}

% ---------------------------------------------------------------
\subsubsection*{2. Zulässige Achsen \texorpdfstring{$A(M_6)$}{A(M_6)}}

$M_6$ führt eine neue Klasse von Achsen ein:

\[
A_{\mathrm{Konzept}} = \{ \text{konzeptionelle Unterscheidungen, die sich aus Attraktoren ergeben} \}.
\]

Diese Achsen entsprechen bedeutungstragenden Freiheitsgraden.

Weitere zulässige Achsen:

\begin{itemize}
    \item $A_{\mathrm{sem}}$ – semantische Einbettungsachsen im neuronalen Raum,
    \item $A_{\mathrm{wm}}$ — Arbeitsspeicher-Erhaltungsachse,
    \item $A_{\mathrm{att}}$ – Aufmerksamkeitsmodulationsachse,
    \item $A_{\mathrm{proto\_logic}}$ – Achsen, die protologische Operationen unterstützen,
    \item $A_{\mathrm{prediction}}$ – prädiktive Codierungsachsen.
\end{itemize}

Eine notwendige Bedingung für $K_6$:

\[
\dim(A_{\mathrm{sem}}) \ge 1
\quad\text{und}\quad
A_{\mathrm{Konzept}} \subseteq A(M_6).
\]

% ---------------------------------------------------------------
\subsubsection*{3. Potentiale \texorpdfstring{$P(M_6)$}{P(M_6)}}

Kognitive Potenziale verallgemeinern $M_5$ elektrodynamische Potenziale auf eine Struktur höherer Ebene:

\[
P(M_6) =
(U_{\mathrm{att}},\
U_{\mathrm{wm}},\
U_{\mathrm{Konzept}},\
U_{\mathrm{pred}},\
U_{\mathrm{Stabilität}}).
\]

Interpretation:

\begin{itemize}
    \item $U_{\mathrm{att}}$ – potenzielle Landschaft, die Aufmerksamkeitszyklen prägt,
    \item $U_{\mathrm{wm}}$ – Energie, die zur Aufrechterhaltung der Arbeitsgedächtniszustände erforderlich ist,
    \item $U_{\mathrm{concept}}$ — konzeptionelles Potenzial, das Attraktoren trennt,
    \item $U_{\mathrm{pred}}$ – Potenzial prädiktiver Fehlanpassungen,
    \item $U_{\mathrm{Stabilität}}$ – globales Stabilisierungspotential für die Attraktorgeometrie.
\end{itemize}

Ein kognitives Kontinuum erfordert:

\[
U_{\mathrm{Konzept}} \text{ hat mehrere Minima, die Konzepten entsprechen}.
\]

% ---------------------------------------------------------------
\subsubsection*{4. Schwellenwerte \texorpdfstring{$\Theta(M_6)$}{\Theta(M_6)}}

Zu den wichtigsten Schwellenwerten gehören:

\begin{itemize}
    \item $\Theta_{\mathrm{att}}$ – minimale Aktivierung für Aufmerksamkeitseinbindung,
    \item $\Theta_{\mathrm{wm}}$ – Schwelle für stabile Speichererhaltung,
    \item $\Theta_{\mathrm{sep}}$ — Trennbarkeitsschwelle zwischen Attraktoren,
    \item $\Theta_{\mathrm{concept}}$ – Schwelle für die Entstehung unterschiedlicher Konzepte,
    \item $\Theta_{\mathrm{coh}}$ – kognitive Kohärenzschwelle.
\end{itemize}

Kritische Voraussetzung für die Existenz eines kognitiven Kontinuums:

\[
\Theta_{\mathrm{sep}} < \Delta_{\mathrm{basin}},
\quad
\Theta_{\mathrm{coh}} \le T_{\mathrm{cog}},
\]
wobei $\Delta_{\mathrm{basin}}$ der Abstand zwischen den Attraktoren ist.

% ---------------------------------------------------------------
\subsubsection*{5. Flüsse \texorpdfstring{$J(M_6)$}{J(M_6)}}

Zu den zulässigen kognitiven Flüssen gehören:

\[
J(M_6) =
\big(
J_{\mathrm{att}},\
J_{\mathrm{wm}},\
J_{\mathrm{Konzept}},\
J_{\mathrm{pred}},\
J_{\mathrm{Plastizität}},\
\nabla_i T_6
\Big).
\]

Interpretation:

\begin{itemize}
    \item $J_{\mathrm{att}}$ – Aufmerksamkeitsmodulationsfluss,
    \item $J_{\mathrm{wm}}$ – Arbeitsspeicher-Stabilisierungsfluss,
    \item $J_{\mathrm{concept}}$ – Strömungen, die die Attraktorgeometrie transformieren,
    \item $J_{\mathrm{pred}}$ – Vorhersagefehlerausbreitung,
    \item $J_{\mathrm{Plastizität}}$ – synaptische Aktualisierungen, die die Kohärenz aufrechterhalten.
\end{itemize}

Erhaltungsbeziehung:

\[
\sum_i J_{\mathrm{Plastizität},i} = \frac{d\mathcal{W}}{dt}.
\]

% ---------------------------------------------------------------
\subsubsection*{6. Zyklen \texorpdfstring{$C(M_6)$}{C(M_6)}}

Schlüsselzyklen, die die Erkenntnis ermöglichen:

\[
C_{\mathrm{att\_wm}}
=
\{\text{Aufmerksamkeit} \to \text{Kodierung} \to
\text{Wartung} \to \text{Aktualisierung} \to \text{Veröffentlichung}\}.
\]

Weitere Zyklen:

\begin{itemize}
    \item $C_{\mathrm{predictive}}$ – Vorhersage $\to$ Fehler $\to$ Update,
    \item $C_{\mathrm{concept}}$ — Stabilisierung konzeptioneller Grenzen,
    \item $C_{\mathrm{proto\_logic}}$ — logische Transformationen von Zuständen,
    \item $C_{\mathrm{Plastizität}}$ – Lernzyklen, die $\mathcal{W}$ anpassen.
\end{itemize}

Notwendige Voraussetzung für das kognitive Leben:

\[
\oint_{C_{\mathrm{att\_wm}}} dA_{\mathrm{Zustand}} \ approx 0.
\]

% ---------------------------------------------------------------
\subsubsection*{7. Grenze \texorpdfstring{$\partial\Omega(M_6)$}{\partial\Omega(M_6)}}

Eine Konfiguration erreicht $\partial\Omega(M_6)$, wenn:

\begin{itemize}
    \item-Attraktoren verlieren an Stabilität oder kollabieren,
    \item semantische Achsen werden degeneriert,
    \item-Vorhersagefehler divergiert (unbegrenzt),
    \item Plastizität bricht die Kohärenz,
    \item $\Theta_{\mathrm{concept}}$ kann nicht gekreuzt werden (keine Konzeptbildung),
    \item Aufmerksamkeit kann sich nicht stabilisieren (Zyklus schlägt fehl).
\end{itemize}

Das Überschreiten der Grenze impliziert den Zusammenbruch von $K_6$ in eine $K_5$-ähnliche Dynamik (rein elektrisch).

% ---------------------------------------------------------------
\subsubsection*{8. Beziehung zu \texorpdfstring{$M_5$}{M_5} und \texorpdfstring{$M_7$}{M_7}}

\textbf{Von $M_5$ bis $M_6$:}

Anforderungen:

\[
C_{\mathrm{spike}} \text{ muss in strukturierte Muster einkoppeln},
\quad
\mathcal{W} \text{ muss Attraktorbecken unterstützen}.
\]

Dadurch erhalten Spike-Züge eine bedeutungstragende Geometrie.

\textbf{Auf dem Weg zu $M_7$ (soziale Kognition):}

\[
A_{\mathrm{shared}} \in A(M_7),
\qquad
\text{erfordert ein stabiles konzeptionelles Repertoire in } M_6,
\]
Sie bilden die Grundlage für Kommunikation, gemeinsame Aufmerksamkeit und soziale Schlussfolgerungen.
